% ===================================================================
\section{Putting it all together}
% ===================================================================

\subsection{The wrap-up lecture's messages}
The final lecture (13 Nov) warns that its slides are ``not your summary sheet for the exams''; its chief messages, with
where each sits in these notes:
\begin{center}\small
\begin{tabularx}{\textwidth}{lX}
\toprule
theme & message\\\midrule
biology & many structural and temporal levels interact and change simultaneously: that causes the complexity (\S2)\\
modelling & useful, but never mistake the model for reality; mind its limitations (\S2, \S4)\\
population numbers & the strategy depends on problem, data and insight needed: no universal strategy (\S3, \S4)\\
complex systems & high-dimensional objects can have simple low-dimensional projections; a model with the expected behaviour
is not thereby appropriate (\S4)\\
applications & done properly, models inform policy and save lives; done badly, they harm (\S3, \S8)\\
competition & competitors often drive others extinct (water hyacinth in the Mula river at Dapodi); coexisting competitors
probably partition the niche, though it is often hard to show (\S5)\\
predation & in simple homogeneous systems prey are driven extinct; real systems are heterogeneous and both sides evolve: a
dynamic, ever-shifting equilibrium (\S6)\\
organism $\leftrightarrow$ environment & a two-way interaction, which is why ecosystems are hard to change (\S7)\\
conservation & very complex; needs inputs from many sources (\S8)\\
evolution & complex; much known, much to learn; work in progress, no major paradigm shift yet (\S9)\\
Hardy--Weinberg & ``the law is boring; the deviations are interesting''; even simple departures give counter-intuitive results (\S10)\\
selection & NS = differential survival and reproduction due to phenotype; NS is not evolution; evolution can be fast;
selection's effects are not always easy to predict (\S11)\\\bottomrule
\end{tabularx}
\end{center}

\subsection{Threads that cross chapters}
\begin{chain}[: one mathematics of thresholds]
\begin{itemize}
\item \textbf{Lakes}: vegetation--turbidity loop $\Rightarrow$ two states, folds, hysteresis.
\item \textbf{Allee effect}: positive density dependence $\Rightarrow$ extinction threshold $A$; harvest $\Rightarrow$ fold.
\item \textbf{Competition, unstable case}: neither species can invade $\Rightarrow$ saddle, priority effect.
\item \textbf{Underdominance}: heterozygote disadvantage $\Rightarrow$ unstable allele frequency.
\item \textbf{Passenger pigeon}: social facilitation $\Rightarrow$ MVP crossed while common.
\item \textbf{Wildebeest}: migration appears only above a tracking radius of $\sim$80\,km.
\end{itemize}
Common lesson: once past the threshold, reversing the original pressure does not restore the old state. Act before it, and
monitor the margin (resilience), not just the current state.
\end{chain}

\begin{chain}[: stability, three ways]
For a continuous model a fixed point is stable if the slope $f'(x^*)<0$ (logistic at $K$: $-r$). For a discrete map, if
$|g'(x^*)|<1$ (Ricker: $|1-r|<1$); losing stability at $g'=-1$ gives period doubling. For a linear matrix model, the zero
state is unstable when the dominant eigenvalue $\lambda>1$. In 2-D (competition, predator--prey) the Jacobian's trace and
determinant decide: trace $<0$ and determinant $>0$ means stable; purely imaginary eigenvalues (basic L--V) mean neutral.
\end{chain}

\begin{chain}[: chance at every level]
Demographic stochasticity (extinction of small populations), genetic drift ($H_t=H_0(1-1/2N_e)^t$), environmental
stochasticity, the integer lattice (chaos to cycles or extinction), founder bottlenecks in invaders, and chance in the
``mess'' of selection. Small numbers make all of them worse; large numbers average out only the demographic and genetic kinds.
\end{chain}

\begin{chain}[: top-down and bottom-up]
Hare--lynx: three trophic levels interact (Krebs). Lakes: nutrients (bottom-up) and fish $\to$ \emph{Daphnia} $\to$ algae
(top-down: biomanipulation, Lake Washington's second jump). Wildebeest: food-limited (bottom-up); zebra probably
predation-limited (top-down). Volterra's principle: indiscriminate killing helps prey.
\end{chain}

\subsection{Past-paper topics outside this year's slides}
Several past papers (2021--2025) asked about topics that are not on this year's slides; they come from the self-study items
of lecture 2 (temperature, biogeochemical cycles) and from earlier runs of the course (ageing). They are summarised here as
general knowledge, in case they reappear.

\textbf{Ageing (evolutionary theories).} Both rest on the \textbf{declining force of selection with age}: few individuals
survive to old age, so effects acting late matter little to fitness.
\begin{itemize}
\item \textbf{Mutation accumulation} (Medawar): late-acting deleterious mutations are weakly selected against and
accumulate by mutation and drift; ageing is their summed effect. Predicts increased genetic variance in late-life traits.
\item \textbf{Antagonistic pleiotropy} (Williams): alleles with early benefits and late costs are actively \emph{favoured};
ageing is the price of early vigour. Predicts negative genetic correlation between early and late fitness.
\item \textbf{Extrinsic mortality} sets the strength of selection on late life: low extrinsic mortality (captivity, flight,
shells) $\Rightarrow$ selection maintains late fitness $\Rightarrow$ slower ageing. So zoo-bred rabbits evolve to age
\emph{slower}; zoo rabbits released to the wild evolve to age \emph{faster}. A deleterious allele acting at 35--55 is purged
less than one acting at 15--35.
\item Birds outlive same-sized mammals despite higher metabolic rates: evidence against the \textbf{rate-of-living} theory.
\item ``Programmed death for the good of the species'' relies on group selection (a cheater allele that keeps its carrier
alive would spread) and is contradicted by the observation that wild animals rarely live long enough to age.
\end{itemize}

\textbf{Nutrient cycles and agriculture.} Energy flows, matter cycles: reservoirs and fluxes; gaseous cycles (C, N) vs
sedimentary (P, S). \textbf{Nitrogen fixation}: N$_2\to$ NH$_3$/NH$_4^+$ by nitrogenase (\emph{Rhizobium} in legume
nodules, cyanobacteria), lightning, industry. \textbf{Nitrification}: NH$_4^+\to$ NO$_2^-\to$ NO$_3^-$ by aerobic
bacteria. \textbf{Denitrification}: NO$_3^-\to$ N$_2$O $\to$ N$_2$ by anaerobic bacteria; returns N to the air.
\textbf{Crop rotation} alternates crops (e.g.\ with N-fixing legumes) on one field; \textbf{slash-and-burn} (\emph{jhum})
releases nutrients as ash, farms a few years, then leaves a long fallow; sustainable only with long fallows and low density.
\textbf{Mean residence time} = stock/throughput flux. Phosphorus, unlike N, has no gaseous phase and is relatively
immobile in soil (hence the past question about advising a king with P-deficient crops: return P via manure, bones, ash,
river silt).

\textbf{Temperature} (Begon \& Townsend, self-study). Rates rise roughly exponentially with temperature: $Q_{10}$, the factor
per 10\,°C, often $\approx2$. Development is summarised by \textbf{day-degrees} above a threshold. The
\textbf{temperature--size rule}: final size falls as rearing temperature rises. \textbf{Endotherms} hold body temperature
nearly constant (flat line against environment) at high energy cost; \textbf{ectotherms} track environmental temperature
(near the line $T_b=T_{env}$), adjusted by basking and shade-seeking; insects can warm up by basking or shivering their
flight muscles. Freezing removes liquid water from cells, so cold injury resembles drought and salinity stress.

\begin{trap}
These items are not in this year's slides; if a question uses them, reason from the principle (declining force of selection;
mass balance of nutrients; rate--temperature relations) rather than recalled detail.
\end{trap}
